%% ===================================================================
\section{Exact levels in families}
\label{sec:synthesis}
%% ===================================================================

By \Cref{thm:dichotomy}, determining the level of a pair
$(G,\mathcal C)$ is a question about subgroups: one has to find the
proper $\mathcal C$-generated subgroups and their derived lengths. This
section does so for symmetric and alternating groups, for the minimal
simple groups, for transvections in $\SL_n(q)$ and for congruence
quotients of $\GL_2$. The involution classes in rank one are treated in
\Cref{sec:rankone}, and \Cref{tab:synth} collects the resulting levels.

\subsection{Symmetric and alternating groups}

\begin{corollary}[Symmetric groups: obstruction at every depth]
\label{cor:Sn}
For $G=S_n$, $n\ge5$, with $\mathcal{C}$ the class of transpositions,
a uniform strict-contraction level exists exactly when $n=5$, and its exact
value is $3$.
For every $n\ge6$ and every depth $m\ge0$, there is a valid configuration $\tau$
with all leaves in the point stabilizer $S_{n-1}<S_n$ whose operator satisfies
\[
 \opn{M_\tau(\rho_{\mathrm{std}})}=1,
\]
where $\rho_{\mathrm{std}}$ is the standard representation of $S_n$.
Consequently no depth gives strict contraction uniformly over all valid
configurations for any $n\ge6$.
\end{corollary}
\begin{proof}
A subgroup generated by transpositions is a Young subgroup $S_{a_1}\times\cdots\times
S_{a_r}$, corresponding to the connected components of the transposition
graph. For $n=5$, every proper such subgroup is solvable of derived length
at most $3$, and the point stabilizer $S_4$ attains $3$; hence
\Cref{thm:dichotomy} gives the exact level.

For $n\ge6$, let $H=S_{n-1}$ fix the last point. This is a proper
non-solvable subgroup generated by transpositions. By \Cref{lem:central},
the targets of configurations with leaves in the transposition class of
$H$ generate $H^{(m)}\ne1$ at every depth $m$. Thus at least one such
configuration is valid at each depth. Every leaf fixes
\[
 v=(1,\ldots,1,-(n-1))\ne0,
 \qquad v\in\{x\in\C^n:\textstyle\sum_i x_i=0\},
\]
so every leaf projection, and therefore the entire operator, fixes $v$.
As the operator is a product of contractions, its standard-representation
norm is exactly $1$. This gives the claimed obstruction at every depth.
\end{proof}

Exact shuffles by freely chosen lazy transpositions exist for every $S_n$
\cite{AngelHolroyd2018}, with suitable probabilities.
\Cref{cor:Sn} identifies when fair independent reads guarantee strict
contraction over all valid commutator trees.

\begin{corollary}[The alternating group $A_5$]
\label{cor:A5}
For $G=A_5$ with $\mathcal{C}$ the class of involutions, the exact level
is $2$.
\end{corollary}
\begin{proof}
The proper $\mathcal{C}$-generated subgroups are $C_2$, $V_4$, $S_3$ and $D_{10}$, all
solvable, with $\max\operatorname{dl}=\operatorname{dl}(S_3)=\operatorname{dl}(D_{10})=2$.
The subgroup $A_4$ is not $\mathcal{C}$-generated, since its involutions generate only
$V_4$. Hence $D=2$ by \Cref{thm:dichotomy}(2).
\end{proof}

\subsection{Minimal simple groups}

\begin{corollary}[Minimal simple groups]
\label{cor:minsimple}
Let $G$ be a finite non-abelian simple group. The following are equivalent:
\begin{enumerate}[nosep,leftmargin=1.6em]
\item $G$ is minimal simple: every proper subgroup is solvable.
\item Every non-trivial conjugacy class has a uniform strict-contraction
level.
\item Every non-trivial conjugacy class has uniform strict contraction at
level $5$.
\end{enumerate}
When these conditions hold, the exact level from \Cref{thm:dichotomy}
satisfies $D(G,\mathcal C)\le5$ for every non-trivial conjugacy class
$\mathcal C$.
The universal bound $5$ is attained by the involution class of
$\PSL_3(3)$ (\Cref{prop:psl33-level}).
\end{corollary}
\begin{proof}
Every non-trivial conjugacy class $\mathcal C$ generates $G$, since
$\langle\mathcal C\rangle$ is non-trivial and normal. If $G$ is minimal
simple, all proper $\mathcal C$-generated subgroups are solvable and have
derived length at most $5$~\cite[Remark 2.3]{ErsoyTortoraTota2014}.
\Cref{thm:dichotomy} gives the exact-level formula and (1)$\Rightarrow$(3).
The implication (3)$\Rightarrow$(2) is immediate.

For (2)$\Rightarrow$(1), suppose that $H<G$ is non-solvable. Choose
$x\in H\setminus O_\infty(H)$, where $O_\infty(H)$ is the largest soluble
normal subgroup of $H$, and set $\mathcal C=x^G$ and
$L=\langle\mathcal C\cap H\rangle$. Conjugation by $H$ preserves
$\mathcal C\cap H$, so $L\trianglelefteq H$. Since $x\in L$, solvability
of $L$ would contradict $x\notin O_\infty(H)$. Moreover $L\le H<G$ and
$\mathcal C\cap L=\mathcal C\cap H$, so $L$ is a proper non-solvable
$\mathcal C$-generated subgroup. \Cref{thm:dichotomy}(1) gives an
obstruction at every depth, contradicting (2).
\end{proof}

Thus contraction for every non-trivial class singles out precisely the
groups in Thompson's classification~\cite{Thompson1968}:
$\PSL_2(2^p)$ ($p$ prime), $\PSL_2(3^p)$ ($p$ odd prime), $\PSL_2(p)$ for
primes $p>3$ with $p\equiv\pm2\pmod5$, the Suzuki groups
$\mathrm{Sz}(2^p)$ ($p$ odd prime), and $\PSL_3(3)$.

\begin{proposition}[The involution class of $\PSL_3(3)$]
\label{prop:psl33-level}
For $G=\PSL_3(3)$ and its involution class, the exact level is $5$.
\end{proposition}
\begin{proof}
Here $G=\SL_3(3)$ has trivial centre. Its involutions have eigenvalues
$1,-1,-1$, form a single conjugacy class, and generate $G$ by simplicity.
The maximal-subgroup list in the ATLAS~\cite{ATLASL33} consists of
two parabolic classes $3^2{:}\GL_2(3)$, one class of $S_4$, and $13{:}3$.

Write a point parabolic as $P=V\rtimes L$, where
$V\cong\F_3^2$ and $L\cong\GL_2(3)$. In the Levi factor,
\[
 L'=\SL_2(3),\qquad \SL_2(3)'=Q_8,\qquad
 Q_8'=\langle z\rangle,\qquad z=-I.
\]
The element $z$ acts as negation on $V$, so $[V,\langle z\rangle]=V$.
Consequently the derived series of $P$ is
\[
 P>V\rtimes\SL_2(3)>V\rtimes Q_8>
 V\rtimes\langle z\rangle>V>1,
\]
and $\operatorname{dl}(P)=5$.

The group $L$ is generated by
$s=\operatorname{diag}(-1,1)$ and the elementary transvections
$t=I+E_{12},I+E_{21}$. Since $sts=t^{-1}$ for either transvection,
both $s$ and $st$ are involutions, so $L$ is involution-generated.
For every $v\in V$, the product $(v z v^{-1})z$ is the translation
$2v$. Thus the involutions of $P$ also generate $V$, and hence $P$.
The line parabolic has the same properties by inverse transpose.

The remaining maximals are $S_4$, of derived length $3$, and
the odd-order group $13{:}3$. Every proper involution-generated
subgroup therefore has derived length at most $5$, and a parabolic
attains $5$.
Apply \Cref{thm:dichotomy}.
\end{proof}

\begin{remark}[$\PSL_2(7)$]
\label{rem:gl32}
The group $\GL_3(\F_2)\cong\PSL_2(7)$ is minimal simple, with maximal
subgroups two classes of $S_4$ and one of $7{:}3$. For the involution
class the proper $\mathcal{C}$-generated subgroups have maximal derived
length $\operatorname{dl}(S_4)=3$, the group $7{:}3$ having odd order,
so $D=3$ by \Cref{thm:dichotomy}(2).
\end{remark}

\subsection{Transvections}
\label{sec:layers}

For $G=\GL_n(\F_q)$ or $\SL_n(\F_q)$ the natural candidates for
obstructions are subspace stabilisers. If $P$ is a proper parabolic
subgroup and the image of $\langle P\cap\mathcal C\rangle$ in a Levi
factor of $P$ is non-solvable, then $\langle P\cap\mathcal C\rangle$ is
a proper non-solvable $\mathcal C$-generated subgroup, and
\Cref{thm:dichotomy}(1) applies. For transvections this settles all but
two cases.

\begin{proposition}[Transvection leaves]
\label{prop:transvection-levels}
Let $n\ge3$, let $G=\SL_n(q)$, and let $\mathcal T$ be its class of
transvections. The exact level is $3$ for
$(n,q)=(3,2)$ and $4$ for $(n,q)=(3,3)$. In every other case there is
a permanent obstruction.
\end{proposition}
\begin{proof}
Transvections generate $\SL_n(q)$. For $q\ge4$ the upper-left
$\SL_2(q)$ is a proper non-solvable $\mathcal T$-generated subgroup.
For $q\in\{2,3\}$ and $n\ge4$, use the upper-left $\SL_3(q)$ instead.
Apply \Cref{thm:dichotomy}(1). In $\SL_3(2)$ transvections are precisely
the involutions, and \Cref{rem:gl32} gives level $3$.

For $\SL_3(3)=\PSL_3(3)$ use the maximal-subgroup list in
\Cref{prop:psl33-level}. In a point parabolic $P=3^2{:}\GL_2(3)$,
a transvection acts trivially on the invariant line, so its Levi image
lies in $\SL_2(3)$. The radical $V\cong\F_3^2$ and the elementary
matrices in this Levi factor are transvections; hence
$\langle P\cap\mathcal T\rangle=J=V\rtimes\SL_2(3)$.
Since $\SL_2(3)'=Q_8$, $Q_8'=\langle-I\rangle$, and $-I$ acts as
negation on $V$, the derived series is
\[
 J\ >\ V\rtimes Q_8\ >\ V\rtimes\langle-I\rangle\ >\ V\ >\ 1,
\]
with orders $216,72,18,9,1$. The dual parabolic has the same property.
The other maximals, $S_4$ and $13{:}3$, have derived lengths $3$ and $2$.
Every proper $\mathcal T$-generated subgroup therefore has derived
length at most $4$, and $J$ attains $4$. Apply \Cref{thm:dichotomy}(2).
\end{proof}

\subsection{Abelian and congruence quotients}

\begin{remark}[One-dimensional representations]
\label{rem:abelian}
Let $\chi$ be a one-dimensional representation of $G$ and let
$w(\tau)=(g_1,\ldots,g_N)$. By \Cref{lem:factor},
\[
 M_\tau(\chi)=\prod_{j=1}^N\tfrac12\bigl(1+\chi(g_j)\bigr),
 \qquad
 \abs{M_\tau(\chi)}=\prod_{j=1}^N\abs{\cos(\theta_j/2)},
\]
where $\chi(g_j)=e^{i\theta_j}$. This has modulus one exactly when
$\chi$ is trivial on every leaf, in accordance with \Cref{lem:geom}. In
particular, if some leaf lies outside the derived subgroup $G'$, the
output law is not supported on a single coset of $G'$. For $G=\GL_n(\F_q)$ with $q>3$, involutions have determinant
$\pm1$ and cannot generate $G$, so a generating leaf set must contain
elements whose determinants generate $\F_q^\times$.
\end{remark}

Let $p$ be a prime and $r\ge2$. Reduction modulo $p$ is a surjection
$\GL_2(\Z/p^r)\to\GL_2(\F_p)$ whose kernel
$K=1+pM_2(\Z/p^r)$ is a $p$-group, so $\GL_2(\Z/p^r)$ is non-solvable
exactly when $p\ge5$. For conjugacy classes, Serre's lifting lemma
reduces the dichotomy to the residue field.

\begin{proposition}[Congruence lifting with generating determinants]
\label{prop:congruence-class}
Let $p\ge5$ be prime, $r\ge2$, $G=\GL_2(\Z/p^r)$, and
$K=\ker(G\to\GL_2(\F_p))$. Let $\mathcal C$ be conjugation-closed,
generate $G$, and have the property that $\det c$ generates
$(\Z/p^r)^\times$ for every $c\in\mathcal C$.
Then $(G,\mathcal C)$ has a uniform contraction level if and only if
$(\GL_2(\F_p),\bar{\mathcal C})$ does. In this case their exact levels
satisfy
\[
 \bar D\le D\le\bar D+\operatorname{dl}(K).
\]
In particular, $\bar D\le D\le\bar D+1$ for $r=2$.
The determinant hypothesis holds for every single conjugacy class
that generates $G$.
\end{proposition}
\begin{proof}
Every non-trivial $\mathcal C$-generated subgroup $H$ has full
determinant image. If its reduction is $\GL_2(\F_p)$, Serre's lifting
lemma~\cite[Lemma 3.15]{Lombardo2015}, applied to the closed preimage of
$H$ in $\GL_2(\Z_p)$, implies that $H$ contains $\SL_2(\Z/p^r)$.
The full determinant image then gives $H=G$. Consequently every proper
$\mathcal C$-generated subgroup has proper reduction. Since $K$ is
solvable, a non-solvable obstruction above yields one below; conversely,
\Cref{prop:quotient} lifts every obstruction below.
In the unobstructed case,
\[
 \operatorname{dl}(H)\le
 \operatorname{dl}(\bar H)+\operatorname{dl}(H\cap K)
 \le\bar D+\operatorname{dl}(K)
\]
for every proper $\mathcal C$-generated $H$. The opposite inequality
$\bar D\le D$ follows from \Cref{prop:quotient}.
For $r=2$, $K\cong M_2(\F_p)^+$ is abelian. Finally, all elements of
a single conjugacy class have the same determinant, which must generate
the unit group if the class generates $G$.
\end{proof}

\begin{remark}[A congruence obstruction for a union of classes]
\label{rem:congruence-supplement}
In $G=\GL_2(\Z/25)$ put $c=\operatorname{diag}(7,1)$ and
$\mathcal C=c^G\cup\{6I\}$. Then
$H=\langle c^G\rangle=\det^{-1}(\langle7\rangle)$ is a proper
non-solvable $\mathcal C$-generated subgroup of index $5$.
Indeed, conjugating $c$ by elementary unipotents and multiplying by
$c^{-1}$ gives all upper and lower unipotents: the respective
coefficients $1-7=-6$ and $1-7^{-1}=8$ are units modulo $25$.
They generate $\SL_2(\Z/25)$, proving the assertion about $H$.
The determinants $7$ and $\det(6I)=11$ have orders $4$ and $5$, so
$\mathcal C$ generates $G$.
Its reduction $\bar{\mathcal C}=\operatorname{diag}(2,1)^{\GL_2(5)}
\cup\{I\}$ is unobstructed. Every non-trivial
$\bar{\mathcal C}$-generated subgroup has full determinant image;
every non-solvable subgroup of $\GL_2(5)$ contains $\SL_2(5)$ by
Dickson's classification~\cite[Theorem 3.13]{Lombardo2015}, so these
two properties force it to be the whole group.
\end{remark}

\begin{remark}[Questions]
\label{rem:families-questions}
Two questions about matrix groups remain open. For $n\ge3$, which
generating conjugacy classes of $\GL_n(\F_q)$ with irreducible
characteristic polynomial have only solvable proper class-generated
subgroups? Such elements stabilise no proper non-zero subspace, so the
question concerns the non-parabolic maximal subgroups, in contrast with
\Cref{prop:transvection-levels}. Secondly, for a generating conjugacy
class of $\GL_2(\Z/p^2)$, $p\ge5$, with unobstructed reduction, when is
the level in \Cref{prop:congruence-class} equal to $\bar D$ and when to
$\bar D+1$?
\end{remark}
